\begin{table}[H]
\raggedright
\caption{Daily dysfunction: within-person day-level models.}
\label{tab:sup-dysfunction}
\footnotesize
\begingroup
\setlength{\tabcolsep}{4pt}
\renewcommand{\arraystretch}{1.08}
\noindent
\begin{tabular*}{\textwidth}{@{\extracolsep{\fill}}llccc@{}}
\toprule
Evening outcome & Day predictor (within-person) & Estimate & SE & $p$ \\
\midrule
pain interference & Pain intensity & -0.193 & 0.096 & .045 \\
pain interference & Attention to pain & -0.015 & 0.097 & .878 \\
pain interference & Threat value & 0.243 & 0.119 & .041 \\
pain interference & Goal engagement & 0.071 & 0.058 & .228 \\
goal success & Pain intensity & -0.024 & 0.117 & .840 \\
goal success & Attention to pain & -0.122 & 0.119 & .307 \\
goal success & Threat value & 0.314 & 0.145 & .031 \\
goal success & Goal engagement & 0.045 & 0.071 & .525 \\
\bottomrule
\end{tabular*}
\par\vspace{3pt}\noindent\parbox{\textwidth}{\footnotesize\textit{Note.} Multilevel models with participant random intercepts. Predictors are person-centered day means. Day-to-day (between-day) variance in pain was small, so day-level effects are weak relative to the momentary dynamics.}
\endgroup
\end{table}
